% ===================================================================
%  THEORY-ONLY PAPER STRATEGY
%  How to publish a simulation/theory paper without experimental data
%  Target: Physical Review A / Physical Review Research / New J. Phys.
% ===================================================================

\documentclass[11pt,a4paper]{article}
\usepackage[utf8]{inputenc}
\usepackage[T1]{fontenc}
\usepackage{amsmath,amssymb}
\usepackage{graphicx}
\usepackage{hyperref}
\usepackage{enumitem}
\usepackage{booktabs}

\begin{document}

\title{\bfseries Theory-Only Publication Strategy\\
       for QESPM Without Experimental Data}
\date{14 July 2026}
\maketitle

\section{The Core Problem}

The current manuscript reads like an experimental proposal: calibration
procedures, ISO standards, five-step protocols, equipment lists.  Without
experimental data, this framing is weak --- reviewers will ask ``where is
the experiment?''  The manuscript must be \textbf{repositioned as a
theoretical contribution} that stands on its own mathematical and conceptual
merits.

\section{What Makes a Theory-Only Paper Publishable}

Strong theory papers without experiments share these characteristics:

\begin{enumerate}[label=(\roman*)]
    \item \textbf{A rigorous mathematical result} --- a theorem, a closed-form
          solution, or an analytical bound that was not previously known.
    \item \textbf{A counter-intuitive physical insight} --- something that
          changes how researchers think about the problem.
    \item \textbf{Comprehensive parameter-space exploration} --- phase
          diagrams, scaling laws, and regime maps derived analytically
          and verified numerically.
    \item \textbf{Clear, falsifiable experimental predictions} --- specific
          numbers that an experimental group can test within 1--2 years.
    \item \textbf{Connection to known puzzles} --- the theory explains or
          illuminates an existing experimental anomaly.
\end{enumerate}

\section{Specific Improvements for This Project}

\subsection{1. Analytical scaling laws (replace numerical spot-checks)}

Instead of ``at $h = 40$\,\textmu m, SNR = 80,'' derive \textbf{universal
scaling relations}:

\begin{itemize}
    \item \textbf{Signal scaling:}
    $|\Delta\omega_x| \propto A \cdot k^2 e^{-kh} \cdot (eE_{\rm bg}/2m\omega_x)$
    \item \textbf{Optimal wavenumber:}
    $k_{\rm opt} = 2/h$ (derived from $d/dk[k^2 e^{-kh}] = 0$)
    \item \textbf{Peak signal:}
    $|\Delta\omega_x|_{\rm max} = (eE_{\rm bg}/2m\omega_x) \cdot (4A/h^2)e^{-2}$
    \item \textbf{Resolution--height trade-off:}
    $\delta x_{\rm lat} \cdot \delta z_{\rm vert} \propto h$ (derived analytically)
    \item \textbf{Effective rank scaling:}
    $r_{\rm eff} \propto 1/h$ (from SVD analysis, fitted numerically with
    analytical justification)
\end{itemize}

These scaling laws are \textbf{the primary deliverable} of a theory paper.
They allow any experimentalist to compute the expected signal for their
specific parameters without re-running the simulation.

\subsection{2. Rigorous Floquet analysis (beyond pseudopotential)}

The current derivation uses the pseudopotential approximation
($\omega_{\rm sec} \approx (\Omega/2)\sqrt{a + q^2/2}$).  For a theory
paper, derive the \textbf{exact Floquet exponent} for the Mathieu equation
with a sinusoidal surface perturbation:

\begin{equation}
\ddot{x} + [a - 2q\cos(2\tau) + \varepsilon\sin(\kappa\tau)]x = 0,
\end{equation}

where $\varepsilon \propto A k^2 e^{-kh}$ is the surface perturbation and
$\kappa = 2k v_{\rm scan}/\Omega$ is the Doppler-shifted spatial frequency
experienced by the moving ion.  Compute the Floquet exponent
$\mu(a,q,\varepsilon,\kappa)$ perturbatively and identify the resonance
conditions where $\mu$ becomes imaginary (parametric instability).

\textbf{This is publishable in itself} --- the Floquet analysis of a Paul
trap with a spatially periodic surface perturbation is, to my knowledge,
a new problem in the Mathieu equation literature.

\subsection{3. Phase diagram in $(h, k, A, {\rm SNR})$ space}

Replace isolated numerical examples with a comprehensive phase diagram:

\begin{itemize}
    \item \textbf{Detectability boundary:} ${\rm SNR}(h,k,A) = 1$ surface
          in $(h,k,A)$ space.
    \item \textbf{Perturbation validity boundary:}
          $|\Delta\omega_x|/\omega_x = 0.05$ surface.
    \item \textbf{Lamb-Dicke validity boundary:}
          $\Delta x_{\rm ion}/\lambda_{\rm feature} = 0.1$ surface.
    \item \textbf{Heating-limited boundary:}
          $\dot{\bar{n}} \cdot \tau_{\rm coh} = 1$ surface.
\end{itemize}

The intersection of these boundaries defines the \textbf{operational regime}
where QESPM is both feasible and well-described by the theory.  Plot this
as a 2D projection (e.g.\ $(h, \lambda)$ plane for fixed $A$) with shaded
regions.

\subsection{4. Falsifiable predictions}

A theory paper must give experimentalists something to test.  Five specific,
falsifiable predictions:

\begin{enumerate}[label=P\arabic*., leftmargin=*]
    \item \textbf{Scaling law test:} For a sinusoidal grating with
          $A = 50$\,nm, $\lambda = 20$\,\textmu m, measure
          $|\Delta\omega_x|$ as a function of $h$ from 10--100\,\textmu m
          and verify the $h^{-2}e^{-2\pi h/\lambda}$ scaling.  Deviation
          from this scaling at small $h$ indicates the onset of Casimir--Polder
          forces or boundary perturbation breakdown.
    
    \item \textbf{Band-pass signature:} Scan the ion across a surface with
          broadband roughness (not a single sinusoid).  The power spectrum
          of $\Delta\omega_x(x)$ should show the characteristic
          $k^4 e^{-2kh}$ shape, peaking at $k = 2/h$.  This is a
          \textbf{smoking-gun signature} of the ion-probe ITF.
    
    \item \textbf{Charge elimination test:} Measure $\Delta\omega_x$ before
          and after in-situ UV irradiation.  The difference isolates the
          removable charge contribution.  The residual should correlate
          with AFM-measured topography.
    
    \item \textbf{Micromotion cross-validation:} Simultaneously measure
          $\Delta\omega_x$ (curvature) and $\beta_x$ (gradient).  For a
          Gaussian surface feature, the ratio
          $\beta_x / (\partial\Delta\omega_x/\partial x)$ should be
          constant --- a self-consistency check.
    
    \item \textbf{Quantum projection noise floor:} With the ion cooled to
          $\bar{n} < 0.1$, the Allan deviation of repeated frequency
          measurements should follow $\sigma_\omega(\tau) \propto 1/\sqrt{\tau}$
          down to the quantum projection noise limit, below which it should
          plateau.  The plateau level is predicted by Eq.~(xx).
\end{enumerate}

\subsection{5. Connection to known experimental anomalies}

Link the theory to existing puzzles in the trapped-ion literature:

\begin{itemize}
    \item \textbf{Anomalous heating:} The $h^{-4}$ scaling of heating rates
          was empirically observed by Turchette et al.\ (2000) and
          Deslauriers et al.\ (2006) but the microscopic mechanism remains
          debated.  The QESPM framework provides a \textbf{spatially resolved}
          probe of the noise source distribution --- if heating-rate maps
          correlate with surface topography, the noise mechanism is geometric
          (patch potentials); if they are uniform, it is intrinsic (Johnson
          noise).
    
    \item \textbf{Patch potential correlation length:} The spatial
          correlation length $\ell_c$ of patch potentials is unknown but
          crucial for ion-trap design.  QESPM can measure $\ell_c$ directly
          by scanning the ion and computing the autocorrelation of
          $\Delta\omega_x(x,y)$.  This is a measurement that \textbf{no
          existing technique can perform}.
\end{itemize}

\section{Recommended Manuscript Restructuring}

\begin{center}
\begin{tabular}{@{}p{3.5cm}ll@{}}
\toprule
Current section & Problem & Repositioned as \\
\midrule
``I propose a novel surface metrology modality...'' & Reads as experimental proposal &
``I develop a theoretical framework for surface potential sensing with a single trapped ion...'' \\
Calibration procedure (5 steps) & Too applied for theory paper & Move to Appendix / SI \\
ISO 25178 classification & Premature without data & Move to Discussion as ``future metrological application'' \\
Numerical spot-checks at specific $(h,k,A)$ & Not systematic &
Universal scaling laws + phase diagrams \\
Two-height ``separation'' claim & Now proven impossible &
\textbf{Theorem:} two-height scanning cannot separate topography from charge (rank-1 proof) \\
\bottomrule
\end{tabular}
\end{center}

\section{Proposed New Structure}

\begin{enumerate}
    \item \textbf{Introduction} (1.5 pages) --- reframe as theoretical
          contribution: ``I develop the theoretical framework for a new
          class of quantum sensor...''
    \item \textbf{Forward model and instrument transfer function} (2 pages)
          --- boundary perturbation theory + ITF derivation (already done)
    \item \textbf{Structure of the inverse problem} (2 pages)
          \begin{itemize}
              \item SVD analysis: effective rank, information content
              \item Resolution matrix and spatially-resolved resolution
              \item \textbf{Theorem:} two-height degeneracy (rank-1 proof)
          \end{itemize}
    \item \textbf{Scaling laws and operational regime} (2 pages)
          \begin{itemize}
              \item Analytical scaling relations derived
              \item Phase diagram in $(h, \lambda, A)$ space
              \item Detectability, validity, and heating boundaries
          \end{itemize}
    \item \textbf{Quantum-limited sensitivity} (1.5 pages)
          \begin{itemize}
              \item QFI for coherent, squeezed, and NOON states
              \item Quantum advantage vs.\ spatial frequency
          \end{itemize}
    \item \textbf{Predictions for experiment} (1.5 pages)
          \begin{itemize}
              \item Five falsifiable predictions (P1--P5 above)
              \item Connection to anomalous heating puzzle
              \item Patch potential correlation length measurement
          \end{itemize}
    \item \textbf{Discussion and outlook} (1 page)
    \item \textbf{Appendix:} Calibration considerations, ISO alignment,
          numerical methods
\end{enumerate}

\section{Journal Targeting}

\begin{center}
\begin{tabular}{@{}llll@{}}
\toprule
Journal & Theory-friendly? & Impact & Fit \\
\midrule
\textbf{Physical Review A} & Yes --- strong theory tradition & IF $\sim 2.9$ & \textbf{Best fit}: atomic physics + rigorous theory \\
\textbf{Physical Review Research} & Yes --- fully open access & New, growing & Good fit: interdisciplinary \\
\textbf{New Journal of Physics} & Yes --- open access & IF $\sim 3.5$ & Good fit: quantum + surface science \\
\textbf{Quantum Sci. Technol.} & Accepts theory & IF $\sim 5.5$ & Requires stronger quantum advantage angle \\
\textbf{Phys. Rev. Applied} & Prefers experimental validation & IF $\sim 4.8$ & Only if predictions are very specific \\
\textbf{Phys. Rev. Letters} & Requires major breakthrough & IF $\sim 8.0$ & Only if Floquet theorem is the main result \\
\bottomrule
\end{tabular}
\end{center}

\textbf{Recommendation:} Target \textbf{Physical Review A} with the
restructured theory paper.  The Floquet analysis of a Mathieu equation
with periodic surface perturbation is a natural fit for PRA's atomic
physics readership.  The quantum Fisher information analysis adds the
quantum metrology angle.  If the two-height degeneracy proof is developed
into a full theorem, consider submitting the core mathematical result as
a \textbf{Physical Review Letters} or rapid communication.

\section{Immediate Implementation Plan}

\begin{enumerate}[label=\textbf{Step \arabic*}, leftmargin=*]
    \item \textbf{Add SVD analysis section} --- already computed
          (figQ2\_svd.pdf), needs text in manuscript.
    \item \textbf{Add QFI section} --- already computed (figQ1\_qfi.pdf),
          needs text.
    \item \textbf{Add two-height degeneracy theorem} --- already proven
          (figQ4\_two\_height\_proof.pdf), needs formal statement.
    \item \textbf{Derive analytical scaling laws} --- extract from existing
          equations and present as numbered scaling relations.
    \item \textbf{Remove calibration procedure from main text} --- move to
          Appendix.
    \item \textbf{Remove ISO 25178 from main text} --- condense to one
          sentence in Discussion.
    \item \textbf{Add ``Predictions for experiment'' section} --- with five
          falsifiable predictions.
    \item \textbf{Reposition abstract and introduction} --- ``theoretical
          framework'' not ``new measurement technique.''
    \item \textbf{Add Floquet analysis} --- if time permits; this is the
          strongest theoretical addition but requires significant new
          derivation.  Can be noted as ``in preparation'' if not ready.
\end{enumerate}

\end{document}
